% - range counting：Ghazi - On the Power of Multiple Anonymous Messages
%     - 构建了一颗range 二叉树，每个节点代表一个range，每个节点都做一次shuffle-DP，任意range可以由这颗二叉树上的O(log)个节点恢复
%     - 如何填dummy message：可以直接用dummy符号，或者2U。
% - median：
%     - 二分：Instance-optimal Mean Estimation Under Differential
%     Privacy


\section{Application}
In this section, we apply our framework for static setting in three queries: frequency estimation, range counting, and median, as defined in Section~\ref{sec:Notation}.
For dynamic setting, we apply our framework in bit counting query. We assume that $\varepsilon = \Theta(1)$ and $\log n = \Theta(\log(1/\delta)) = \Theta(\log M)= \Theta(\log T)$.



\subsection{Static Setting}
% please check: which framework should be used to apply to this 3 queries? multiple or single round. 
% can write both


\paragraph{Frequency estimation}
% 1) run B times sum estimation
% 2) qiyao's fe1
% pure-dp use 1; unless use 2
% we do not care pure/approximate DP, so fe1


For the frequency estimation problem, the SOTA solution LWY \cite{luo2022frequency} sends $O(1)$ messages and achieves an $\ell_{\infty}$-error of
$O\big(\tfrac{1}{\varepsilon}\sqrt{\log \tfrac{U}{\beta} \log \tfrac{1}{\delta}}\big)$.
By integrating it into our two-round framework, we have:
\begin{itemize}
% please check the bias
% read the other paper on how to represent the bias
\item With high probability $1-\beta$, the total error is bounded by: 
$$
% O\Big(\frac{m_{\max}(D)}{\varepsilon}\cdot \big(\log\frac{\log M}{\beta}+\sqrt{\log \frac{U}{\beta}\log\frac{m_{\max}(D)}{\delta}}\big)\Big).
O\Big(\frac{m_{\max}(D)}{\varepsilon}\cdot \big(\log\frac{\log M}{\beta}+\sqrt{\log \frac{U}{\beta}\log\frac{1}{\delta}}\big)\Big).
$$
\item The expected number of messages is $O(1)$.
\end{itemize}

By integrating it into our one-round framework, we have:

\begin{itemize}
\item With high probability $1-\beta$, the total error is bounded by: 
$$
% O\Big(\frac{m_{\max}(D)}{\varepsilon}\big(\log \frac{1}{\beta}+\log M \cdot \sqrt{
% \log \frac{U\cdot\log M }{\beta}
% \cdot \log \frac{\log M \cdot m_{\max}(D)}{\delta} }\big)\Big).
O\Big(\frac{m_{\max}(D)}{\varepsilon}\big(\log \frac{1}{\beta}+\log M \cdot \sqrt{
\log \frac{U\cdot\log M }{\beta}
\cdot \log \frac{1}{\delta} }\big)\Big).
$$
% $$O\Big(\frac{m_{\max}(D)}{\varepsilon}\big(\log \frac{\log M}{\beta}+\log M \cdot \sqrt{
% \log \frac{U\cdot\log M }{\beta}
% \cdot \log \frac{\log M \cdot m_{\max}(D)}{\delta} }\big)\Big).$$

\item The expected number of messages is $O(\log M)$.
\end{itemize}


\paragraph{Range counting}
% on the power of Ghazi'20
% rm2 better, which is an advanced version of <on the power ...>

Range counting can be reduced to frequency estimation by building a binary tree over the domain $[0,U]$, where each node represents a range and each level forms a frequency estimation problem.
The SOTA shuffle-DP protocol LYD+ \cite{luo2025rm2} achieves an $\ell_{\infty}$-error of
$O\big(\tfrac{\log^{1.5} U}{\varepsilon} \sqrt{\log \tfrac{U}{\beta} \log \tfrac{1}{\delta}}\big)$, with $O(\log U)$ messages per user.
By integrating it into our two-round framework, we have:
\begin{itemize}
% please check the bias
% read the other paper on how to represent the bias
\item With high probability $1-\beta$, the total error is bounded by: 
$$
% O\Big(\frac{m_{\max}(D)}{\varepsilon}\cdot \big(\log\frac{\log M}{\beta}+\log^{1.5}U\cdot\sqrt{\log \frac{U}{\beta}\log\frac{m_{\max}(D)}{\delta}}\big)\Big).
O\Big(\frac{m_{\max}(D)}{\varepsilon}\cdot \big(\log\frac{\log M}{\beta}+\log^{1.5}U\cdot\sqrt{\log \frac{U}{\beta}\log\frac{1}{\delta}}\big)\Big).
$$
\item The expected number of messages is $O(\log U)$.
\end{itemize}

By integrating it into our one-round framework, we have:
% 可以把多个项分开，我们的framework只引入了一个更小的可以被dominate的term
\begin{itemize}
\item With high probability $1-\beta$, the total error is bounded by: 
% \[
% \begin{aligned}
% &O\Bigg(
% \frac{m_{\max}(D)}{\varepsilon}\Big(
% \log \frac{\log M}{\beta} + \log^{1.5} U \cdot \log M \cdot \\
% &\sqrt{
% \log \frac{U\cdot\log M}{\beta}
% \cdot \log \frac{\log M \cdot m_{\max}(D)}{\delta}
% }
% \Big)
% \Bigg)
% \end{aligned}
% \]
\[
\begin{aligned}
O\Bigg(
\frac{m_{\max}(D)}{\varepsilon}\Big(
\log \frac{1}{\beta} + \log^{1.5} U \cdot \log M \cdot \sqrt{
\log \frac{U\cdot\log M}{\beta}
\cdot \log \frac{1}{\delta}
}
\Big)
\Bigg)
\end{aligned}
\]
% $$O\Big(\frac{m_{\max}(D)}{\varepsilon}\big(\log \frac{\log M}{\beta}+\log M \cdot \sqrt{
% \log \frac{U\cdot\log M }{\beta}
% \cdot \log \frac{\log M \cdot m_{\max}(D)}{\delta} }\big)\Big).$$

\item The expected number of messages is $O(\log U \log M)$.
\end{itemize}


\paragraph{Median}
% why median not mean
% instance-optimal mean estimation under differential privacy. Ziyue
% use <on the power of ...> to execute frequency estimation, in a binary structure.
% but multiple round

Median can be reduced to range counting by finding a value $v$ such that the number of elements smaller than $v$ is approximately half of the total. The SOTA shuffle-DP protocol HLY~\cite{huang2021instance} achieves a rank error of $\frac{\log^{2.5} U}{\varepsilon} \sqrt{\log \frac{U}{\beta}\log \frac{1}{\delta}}$, with $O(\log ^5 U \log \frac{1}{\delta})$.
% here we assume O(log U) = O(log 1/\delta)
By integrating it into our two-round framework, we have:
\begin{itemize}
\item With high probability $1-\beta$, the rank error is bounded by: $$O\Big(\frac{m_{\max}(D)}{\varepsilon}\big( \log \frac{\log M}{\beta}+\log^3 U \sqrt{\log \frac{U\log M}{\beta}\log \frac{1}{\delta}}\big)\Big).$$
\item The expected number of messages is $O(\log ^5 U \log \frac{1}{\delta})$.
\end{itemize}

% use rank error or inf error?
% log times range counting

\subsection{Dynamic Setting}
In the dynamic setting, we focus on the bit counting query.
The SOTA shuffle-DP protocol for bit counting in the dynamic setting TKMS~\cite{pmlr-v202-tenenbaum23a} achieves an error 
% $O(\frac{\sqrt{\log (1/\beta)}\log^{1.5}T+\log^2T}{\varepsilon}\sqrt{\log \frac{1}{\delta}})$
$O(\frac{\log^{1.5}T}{\varepsilon}\sqrt{\log \frac{1}{\delta}})$
with $O(t\log T)$ messages\footnote{Unlike the general dynamic setting, TKMS assumes a public user-to-time mapping where user $i$ contributes exactly once at time $t=i$. 
This is easily adapted to general setting: let all users participate at every time step $t\in[T]$, with user $i$ reporting their actual bit when $t=i$ and a dummy symbol $\perp$ otherwise.}.
%check the (\log 1/\beta) term
By integrating it into our framework, for all $t\in \{1,2,\dots,T\}$ simultaneously, we have:
\begin{itemize}
    \item With high probability $1-\beta$, the total error is bounded by: 
    \[
    O\Big(
    \frac{\log ^{1.5} T \cdot \log M \cdot m_{\max}(D^t)}{\varepsilon} \cdot (\log \frac{T \cdot \log M}{\beta}+\sqrt{\log\frac{1}{\delta}})
    \Big)
    \]
    \item The expected number of messages is $O( t\cdot M\log T)$.
\end{itemize}

